\section{Normas sociales y emociones en sistemas multiagente}

Diversas emociones sociales influyen directamente en el comportamiento social humano. Emociones como culpa, vergüenza, empatía o aprobación social influyen directamente sobre la cooperación, el cumplimiento de normas y la toma de decisiones colectivas.

En sistemas multiagente, diferentes investigaciones han explorado la integración de emociones sociales dentro de arquitecturas normativas con el objetivo de modelar comportamientos más flexibles y socialmente consistentes \cite{argente2022normative}.

Los agentes normativos emocionales combinan mecanismos tradicionales de razonamiento con estados afectivos capaces de modificar el comportamiento del agente en función del contexto social. De esta forma, las emociones actúan como mecanismos reguladores que afectan tanto al cumplimiento de normas como a las interacciones entre agentes.

Por ejemplo, emociones negativas como culpa o vergüenza pueden aumentar la probabilidad de que un agente respete determinadas normas sociales, mientras que emociones positivas como orgullo o aprobación social pueden reforzar comportamientos cooperativos.

\begin{table}[H]
\centering
\caption{Ejemplos de emociones sociales en agentes normativos}
\begin{tabular}{p{4cm} p{8cm}}
\toprule
\textbf{Emoción} & \textbf{Efecto sobre el agente} \\
\midrule
Culpa & Incrementa cumplimiento normativo \\
Vergüenza & Reduce conductas socialmente negativas \\
Orgullo & Refuerza cooperación y reputación \\
Empatía & Favorece colaboración entre agentes \\
\bottomrule
\end{tabular}
\end{table}

La incorporación de emociones sociales resulta especialmente relevante en aplicaciones relacionadas con simulación social, negociación automática, robótica social y sistemas colaborativos.

Sin embargo, la modelización de normas sociales y emociones continúa siendo un problema complejo debido a la fuerte dependencia contextual y cultural de este tipo de comportamientos.